\scanpage{62}
in more detail.

3. \handwritten{\uline}{The} \handwritten{\uline}{subgroup} $\handwritten{\underline}{\handwritten{\Gamma}}_{\handwritten{0}}(q)$ \handwritten{\uline}{for} $q$ \handwritten{\uline}{prime}.

The theory of \handwritten{\S}2 carries through with some minor changes,and the replacement of the operator $\handwritten{\underline}{T_q}$ by the operator $\handwritten{\underline}{U_q}$ defined by
\[
 U_q F(z)=q^{-1}\mathop{\handwritten{\sum}}\limits_{\handwritten{\underline}{\handwritten{\lambda}}=\handwritten{0}}^{\handwritten{q-1}} \handwritten{F\left(\frac{z+\underline{\lambda}}{q}\right)}=\mathop{\handwritten{\sum}}\limits_{\handwritten{n}}a(nq)\,x^n\,.
\]
(M1) need only hold for $V \handwritten{\in}\handwritten{\underline}{\handwritten{\Gamma}}_{\handwritten{0}}(q)$,and in (M2) the \handwritten{\uline}{two} cusps $\handwritten{\underline}{B}$ can be taken as $z=0$ and $z=i\handwritten{\infty}$. The existence of another cusp enables us to have\inserted{entire} functions which are as it were ``cusp'' functions with a zero at $\handwritten{\underline}{z}=i\handwritten{\infty}$ by virtue of having a pole at $\handwritten{\underline}{z}=0$ . With regard to the classical theory it is convenient to discard the cuspforms $\handwritten{\underline}{G(z)}$ and $\handwritten{\underline}{G(qz)}$ on $\handwritten{\underline}{\handwritten{\Gamma}}_{\handwritten{0}}(\handwritten{\underline}{q})$ where $\handwritten{\underline}{G(z)}$ is a cuspform of the same dimension on the full modular group(and so a fortiori on $\handwritten{\underline}{\handwritten{\Gamma}}_{\handwritten{0}}(q)$ ),and to consider $C^0$ as the space of ``newforms'' which are the remaining cuspforms on $\handwritten{\underline}{\handwritten{\Gamma}}_{\handwritten{0}}(\handwritten{\underline}{q})$ of dimension $\handwritten{\underline}{k}$. We then have

\deleted{Theorem}
\begingroup\csname tagsleft@true\endcsname
\[
 \handwritten{\underline}{T_p\text{ and }U_q\text{ map }C\text{ into }C\text{ and }C^{\handwritten{0}}\text{ into }C^{\handwritten{0}}}\,, \tag*{(H1q)}
\]
\[
 \handwritten{\underline}{\text{All the } \deleted{T\unsure{s|p}}\handwritten{T_p}\text{ commute,and commute with }U}_{q}\,, \tag*{(H2q)}
\]
\[
 \handwritten{\underline}{T_p\text{ is an hermitian operator on }C^{\handwritten{0}}\text{ (note }U_q\text{ is not)}}\,, \tag*{(H3q)}
\]
\[
 \handwritten{\underline}{C^{\handwritten{0}}\text{ is finite dimensional}}. \tag*{(H4q)}
\]
\endgroup

This leads (with a little extra work to deal with $U_q$) to

\handwritten{\uline}{Theorem} 2. \handwritten{\uline}{The space of new cuspforms on} $\handwritten{\underline}{\handwritten{\Gamma}}_{\handwritten{0}}(q)$ \handwritten{\uline}{has a basis} $F_i(z)$ \handwritten{\uline}{of forms} \handwritten{\uline}{which are simultaneous eigenforms} of \handwritten{\uline}{all the operators} $T_p$ \handwritten{\uline}{and of} $U_q$ . In \handwritten{\uline}{terms of the Fourier series} $\mathop{\handwritten{\sum}}\limits_{\handwritten{n=1}}^{\handwritten{\infty}}a(n)\,x^n$ of \handwritten{\uline}{a typica}\deleted{\unsure{o|l}}l $F_i(z)$ \handwritten{\uline}{normalized by} $a(1)=1$ this \handwritten{\uline}{is equivalent to the multiplicative} \handwritten{\uline}{relations}
\[
 \begin{aligned}
 a(np)-a(n)a(p)+p^{-k-1}a(n/p)&=0,\qquad (\text{\handwritten{\uline}{all prime} }p\handwritten{\ne}q,\text{\handwritten{\uline}{all integral} } n),\\
 a(nq)-a(n)a(q)&=0. 
 \end{aligned}
\]
\margin{\handwritten{\#}}
(In fact we have,in the case of negative $k$ where Theorem 2 is non-trivial,that $a(q)=\mathbin{\handwritten{\pm}}q^{-k/2\,-1}$. The theory on $\handwritten{\underline}{\handwritten{\Gamma}}_{\handwritten{0}}(m)$ for a general $\handwritten{\underline}{m}$ will appear in a paper by Joseph Lehner and the author). We now state\deleted{\unsure{a form|form}}\inserted{an extension} of Conjecture 1 which is more closely analogous to Theorem 2.
% archive sheet number: 61
